%----------------------------------------------------------------------------------------
%	PACKAGES AND OTHER DOCUMENT CONFIGURATIONS
%----------------------------------------------------------------------------------------

\documentclass[11pt,fleqn]{book} % Default font size and left-justified equations

\usepackage[top=3cm,bottom=3cm,left=3.2cm,right=3.2cm,headsep=10pt,paperwidth=8.5in,paperheight=11in]{geometry} % Page margins

\usepackage{xcolor} % Required for specifying colors by name

\definecolor{mycolor}{RGB}{24,97,243} % Define the orange color used for highlighting throughout the book
\usepackage{color}
\definecolor{mygreen}{rgb}{0,0.6,0}
\definecolor{mygray}{rgb}{0.5,0.5,0.5}
\definecolor{mymauve}{rgb}{0.58,0,0.82}
\definecolor{greencodeblocks}{RGB}{0,160,0}
\definecolor{graycodeblocks}{RGB}{160,160,160}
\definecolor{bluecodeblocks}{RGB}{0,0,255}

% Font Settings
\usepackage{avant} % Use the Avantgarde font for headings
%\usepackage{times} % Use the Times font for headings
\usepackage{mathptmx} % Use the Adobe Times Roman as the default text font together with math symbols from the Sym­bol, Chancery and Com­puter Modern fonts

\usepackage{microtype} % Slightly tweak font spacing for aesthetics
\usepackage[utf8]{inputenc} % Required for including letters with accents
\usepackage[T1]{fontenc} % Use 8-bit encoding that has 256 glyphs

\usepackage{listings} % Para estrofas de codigo
\usepackage{caption} % Para captar label y caption del lstlisting o listing

%%%%%%%%%%%%%%%%%%%%%%%
\definecolor{gray97}{gray}{.97}
\definecolor{gray75}{gray}{.75}
\definecolor{gray45}{gray}{.45}
 
\usepackage{listings}
\lstset{ frame=Ltb,
     framerule=0pt,
     aboveskip=0.5cm,
     framextopmargin=3pt,
     framexbottommargin=3pt,
     framexleftmargin=0.4cm,
     framesep=0pt,
     rulesep=.4pt,
     backgroundcolor=\color{mycolor!5},
     rulesepcolor=\color{black},
     %
     stringstyle=\ttfamily,
     showstringspaces = false,
     basicstyle=\small\ttfamily,
     commentstyle=\color{graycodeblocks},
     keywordstyle=\bfseries\color{greencodeblocks},
     %identifierstyle=\color{blue},
     stringstyle=\color{bluecodeblocks},
     %
     numbers=left,
     numbersep=15pt,
     numberstyle=\small\color{mycolor},
     numberfirstline = false,
     breaklines=true,
     %
     texcl=true, %para abilitar las tildes en el comentario
   }

% minimizar fragmentado de listados
\lstnewenvironment{listing}[1][]
   {\lstset{#1}\pagebreak[0]}{\pagebreak[0]}
 
\lstdefinestyle{consola}
   {basicstyle=\scriptsize\bf\ttfamily,
    backgroundcolor=\color{gray75},
   }
 
\lstdefinestyle{Cpp}
   {language=C++,
   }

%%%%%%%%%%%%%Aumentar sub sub section 

%%%%%%%%%%%%%%%%%
\newcommand{\tab}[1]{\hspace{.2\textwidth}\rlap{#1}} % nuevo comando para tab :)
%%%%%%%%%%%%%%%%%

%%%%%%%%%%%%%%%%%%%%%%%%%%%%5

% Bibliography
\usepackage[style=alphabetic,sorting=nyt,sortcites=true,autopunct=true,babel=hyphen,hyperref=true,abbreviate=false,backref=true,backend=biber]{biblatex}
\addbibresource{bibliography.bib} % BibTeX bibliography file
\defbibheading{bibempty}{}

% Index
\usepackage{calc} % For simpler calculation - used for spacing the index letter headings correctly
\usepackage{makeidx} % Required to make an index
\makeindex % Tells LaTeX to create the files required for indexing

%----------------------------------------------------------------------------------------

%----------------------------------------------------------------------------------------
%	VARIOUS REQUIRED PACKAGES
%----------------------------------------------------------------------------------------

\usepackage{titlesec} % Allows customization of titles

\usepackage{graphicx} % Required for including pictures
\graphicspath{{Pictures/}} % Specifies the directory where pictures are stored

\usepackage{lipsum} % Inserts dummy text

\usepackage{tikz} % Required for drawing custom shapes

\usepackage[english]{babel} % English language/hyphenation

\usepackage{enumitem} % Customize lists
\setlist{nolistsep} % Reduce spacing between bullet points and numbered lists

\usepackage{booktabs} % Required for nicer horizontal rules in tables

\usepackage{eso-pic} % Required for specifying an image background in the title page

%----------------------------------------------------------------------------------------
%	MAIN TABLE OF CONTENTS
%----------------------------------------------------------------------------------------

\usepackage{titletoc} % Required for manipulating the table of contents

\contentsmargin{0cm} % Removes the default margin
% Chapter text styling
\titlecontents{chapter}[1.25cm] % Indentation
{\addvspace{15pt}\large\sffamily\bfseries} % Spacing and font options for chapters
{\color{mycolor!60}\contentslabel[\Large\thecontentslabel]{1.25cm}\color{mycolor}} % Chapter number
{}  
{\color{mycolor!60}\normalsize\sffamily\bfseries\;\titlerule*[.5pc]{.}\;\thecontentspage} % Page number
% Section text styling
\titlecontents{section}[1.25cm] % Indentation
{\addvspace{5pt}\sffamily\bfseries} % Spacing and font options for sections
{\contentslabel[\thecontentslabel]{1.25cm}} % Section number
{}
{\sffamily\hfill\color{black}\thecontentspage} % Page number
[]
% Subsection text styling
\titlecontents{subsection}[1.25cm] % Indentation
{\addvspace{1pt}\sffamily\small} % Spacing and font options for subsections
{\contentslabel[\thecontentslabel]{1.25cm}} % Subsection number
{}
{\sffamily\;\titlerule*[.5pc]{.}\;\thecontentspage} % Page number
[] 

%----------------------------------------------------------------------------------------
%	MINI TABLE OF CONTENTS IN CHAPTER HEADS
%----------------------------------------------------------------------------------------

% Section text styling
\titlecontents{lsection}[0em] % Indendating
{\footnotesize\sffamily} % Font settings
{}
{}
{}

% Subsection text styling
\titlecontents{lsubsection}[.5em] % Indentation
{\normalfont\footnotesize\sffamily} % Font settings
{}
{}
{}
 
%----------------------------------------------------------------------------------------
%	PAGE HEADERS
%----------------------------------------------------------------------------------------

\usepackage{fancyhdr} % Required for header and footer configuration

\pagestyle{fancy}
\renewcommand{\chaptermark}[1]{\markboth{\sffamily\normalsize\bfseries\chaptername\ \thechapter.\ #1}{}} % Chapter text font settings
\renewcommand{\sectionmark}[1]{\markright{\sffamily\normalsize\thesection\hspace{5pt}#1}{}} % Section text font settings
\fancyhf{} \fancyhead[LE,RO]{\sffamily\normalsize\thepage} % Font setting for the page number in the header
\fancyhead[LO]{\rightmark} % Print the nearest section name on the left side of odd pages
\fancyhead[RE]{\leftmark} % Print the current chapter name on the right side of even pages
\renewcommand{\headrulewidth}{0.5pt} % Width of the rule under the header
\addtolength{\headheight}{2.5pt} % Increase the spacing around the header slightly
\renewcommand{\footrulewidth}{0pt} % Removes the rule in the footer
\fancypagestyle{plain}{\fancyhead{}\renewcommand{\headrulewidth}{0pt}} % Style for when a plain pagestyle is specified

% Removes the header from odd empty pages at the end of chapters
\makeatletter
\renewcommand{\cleardoublepage}{
\clearpage\ifodd\c@page\else
\hbox{}
\vspace*{\fill}
\thispagestyle{empty}
\newpage
\fi}

%----------------------------------------------------------------------------------------
%	THEOREM STYLES
%----------------------------------------------------------------------------------------

\usepackage{amsmath,amsfonts,amssymb,amsthm} % For math equations, theorems, symbols, etc

\newcommand{\intoo}[2]{\mathopen{]}#1\,;#2\mathclose{[}}
\newcommand{\ud}{\mathop{\mathrm{{}d}}\mathopen{}}
\newcommand{\intff}[2]{\mathopen{[}#1\,;#2\mathclose{]}}
\newtheorem{notation}{Notation}[chapter]

%%%%%%%%%%%%%%%%%%%%%%%%%%%%%%%%%%%%%%%%%%%%%%%%%%%%%%%%%%%%%%%%%%%%%%%%%%%
%%%%%%%%%%%%%%%%%%%% dedicated to boxed/framed environements %%%%%%%%%%%%%%
%%%%%%%%%%%%%%%%%%%%%%%%%%%%%%%%%%%%%%%%%%%%%%%%%%%%%%%%%%%%%%%%%%%%%%%%%%%
\newtheoremstyle{ocrenumbox}% % Theorem style name
{0pt}% Space above
{0pt}% Space below
{\normalfont}% % Body font
{}% Indent amount
{\small\bf\sffamily\color{mycolor}}% % Theorem head font
{\;}% Punctuation after theorem head
{0.25em}% Space after theorem head
{\small\sffamily\color{mycolor}\thmname{#1}\nobreakspace\thmnumber{\@ifnotempty{#1}{}\@upn{#2}}% Theorem text (e.g. Theorem 2.1)
\thmnote{\nobreakspace\the\thm@notefont\sffamily\bfseries\color{black}---\nobreakspace#3.}} % Optional theorem note
\renewcommand{\qedsymbol}{$\blacksquare$}% Optional qed square

\newtheoremstyle{blacknumex}% Theorem style name
{5pt}% Space above
{5pt}% Space below
{\normalfont}% Body font
{} % Indent amount
{\small\bf\sffamily}% Theorem head font
{\;}% Punctuation after theorem head
{0.25em}% Space after theorem head
{\small\sffamily{\tiny\ensuremath{\blacksquare}}\nobreakspace\thmname{#1}\nobreakspace\thmnumber{\@ifnotempty{#1}{}\@upn{#2}}% Theorem text (e.g. Theorem 2.1)
\thmnote{\nobreakspace\the\thm@notefont\sffamily\bfseries---\nobreakspace#3.}}% Optional theorem note

\newtheoremstyle{blacknumbox} % Theorem style name
{0pt}% Space above
{0pt}% Space below
{\normalfont}% Body font
{}% Indent amount
{\small\bf\sffamily}% Theorem head font
{\;}% Punctuation after theorem head
{0.25em}% Space after theorem head
{\small\sffamily\thmname{#1}\nobreakspace\thmnumber{\@ifnotempty{#1}{}\@upn{#2}}% Theorem text (e.g. Theorem 2.1)
\thmnote{\nobreakspace\the\thm@notefont\sffamily\bfseries---\nobreakspace#3.}}% Optional theorem note

%%%%%%%%%%%%%%%%%%%%%%%%%%%%%%%%%%%%%%%%%%%%%%%%%%%%%%%%%%%%%%%%%%%%%%%%%%%
%%%%%%%%%%%%% dedicated to non-boxed/non-framed environements %%%%%%%%%%%%%
%%%%%%%%%%%%%%%%%%%%%%%%%%%%%%%%%%%%%%%%%%%%%%%%%%%%%%%%%%%%%%%%%%%%%%%%%%%
\newtheoremstyle{ocrenum}% % Theorem style name
{5pt}% Space above
{5pt}% Space below
{\normalfont}% % Body font
{}% Indent amount
{\small\bf\sffamily\color{mycolor}}% % Theorem head font
{\;}% Punctuation after theorem head
{0.25em}% Space after theorem head
{\small\sffamily\color{mycolor}\thmname{#1}\nobreakspace\thmnumber{\@ifnotempty{#1}{}\@upn{#2}}% Theorem text (e.g. Theorem 2.1)
\thmnote{\nobreakspace\the\thm@notefont\sffamily\bfseries\color{black}---\nobreakspace#3.}} % Optional theorem note
\renewcommand{\qedsymbol}{$\blacksquare$}% Optional qed square
\makeatother

% Defines the theorem text style for each type of theorem to one of the three styles above
\newcounter{dummy} 
\numberwithin{dummy}{section}
\theoremstyle{ocrenumbox}
\newtheorem{theoremeT}[dummy]{Theorem}
\newtheorem{problem}{Problem}[chapter]
\newtheorem{exerciseT}{Exercise}[chapter]
\theoremstyle{blacknumex}
\newtheorem{exampleT}{Example}[chapter]
\theoremstyle{blacknumbox}
\newtheorem{vocabulary}{Vocabulary}[chapter]
\newtheorem{definitionT}{Definition}[section]
\newtheorem{corollaryT}[dummy]{Corollary}
\theoremstyle{ocrenum}
\newtheorem{proposition}[dummy]{Proposition}

%----------------------------------------------------------------------------------------
%	DEFINITION OF COLORED BOXES
%----------------------------------------------------------------------------------------

\RequirePackage[framemethod=default]{mdframed} % Required for creating the theorem, definition, exercise and corollary boxes

% Theorem box
\newmdenv[skipabove=7pt,
skipbelow=7pt,
backgroundcolor=black!5,
linecolor=mycolor,
innerleftmargin=5pt,
innerrightmargin=5pt,
innertopmargin=5pt,
leftmargin=0cm,
rightmargin=0cm,
innerbottommargin=5pt]{tBox}

% Exercise box	  
\newmdenv[skipabove=7pt,
skipbelow=7pt,
rightline=false,
leftline=true,
topline=false,
bottomline=false,
backgroundcolor=mycolor!10,
linecolor=mycolor,
innerleftmargin=5pt,
innerrightmargin=5pt,
innertopmargin=5pt,
innerbottommargin=5pt,
leftmargin=0cm,
rightmargin=0cm,
linewidth=4pt]{eBox}	

% Definition box
\newmdenv[skipabove=7pt,
skipbelow=7pt,
rightline=false,
leftline=true,
topline=false,
bottomline=false,
linecolor=mycolor,
innerleftmargin=5pt,
innerrightmargin=5pt,
innertopmargin=0pt,
leftmargin=0cm,
rightmargin=0cm,
linewidth=4pt,
innerbottommargin=0pt]{dBox}	

% Corollary box
\newmdenv[skipabove=7pt,
skipbelow=7pt,
rightline=false,
leftline=true,
topline=false,
bottomline=false,
linecolor=gray,
backgroundcolor=black!5,
innerleftmargin=5pt,
innerrightmargin=5pt,
innertopmargin=5pt,
leftmargin=0cm,
rightmargin=0cm,
linewidth=4pt,
innerbottommargin=5pt]{cBox}

% Creates an environment for each type of theorem and assigns it a theorem text style from the "Theorem Styles" section above and a colored box from above
\newenvironment{theorem}{\begin{tBox}\begin{theoremeT}}{\end{theoremeT}\end{tBox}}
\newenvironment{exercise}{\begin{eBox}\begin{exerciseT}}{\hfill{\color{mycolor}\tiny\ensuremath{\blacksquare}}\end{exerciseT}\end{eBox}}				  
\newenvironment{definition}{\begin{dBox}\begin{definitionT}}{\end{definitionT}\end{dBox}}	
\newenvironment{example}{\begin{exampleT}}{\hfill{\tiny\ensuremath{\blacksquare}}\end{exampleT}}		
\newenvironment{corollary}{\begin{cBox}\begin{corollaryT}}{\end{corollaryT}\end{cBox}}	

%----------------------------------------------------------------------------------------
%	REMARK ENVIRONMENT
%----------------------------------------------------------------------------------------

\newenvironment{remark}{\par\vspace{10pt}\small % Vertical white space above the remark and smaller font size
\begin{list}{}{
\leftmargin=35pt % Indentation on the left
\rightmargin=25pt}\item\ignorespaces % Indentation on the right
\makebox[-2.5pt]{\begin{tikzpicture}[overlay]
\node[draw=mycolor!60,line width=1pt,circle,fill=mycolor!25,font=\sffamily\bfseries,inner sep=2pt,outer sep=0pt] at (-15pt,0pt){\textcolor{mycolor}{R}};\end{tikzpicture}} % Orange R in a circle
\advance\baselineskip -1pt}{\end{list}\vskip5pt} % Tighter line spacing and white space after remark

%----------------------------------------------------------------------------------------
%	SECTION NUMBERING IN THE MARGIN
%----------------------------------------------------------------------------------------

\makeatletter
\renewcommand{\@seccntformat}[1]{\llap{\textcolor{mycolor}{\csname the#1\endcsname}\hspace{1em}}}                    
\renewcommand{\section}{\@startsection{section}{1}{\z@}
{-4ex \@plus -1ex \@minus -.4ex}
{1ex \@plus.2ex }
{\normalfont\large\sffamily\bfseries}}
\renewcommand{\subsection}{\@startsection {subsection}{2}{\z@}
{-3ex \@plus -0.1ex \@minus -.4ex}
{0.5ex \@plus.2ex }
{\normalfont\sffamily\bfseries}}
\renewcommand{\subsubsection}{\@startsection {subsubsection}{3}{\z@}
{-2ex \@plus -0.1ex \@minus -.2ex}
{.2ex \@plus.2ex }
{\normalfont\small\sffamily\bfseries}}                        
\renewcommand\paragraph{\@startsection{paragraph}{4}{\z@}
{-2ex \@plus-.2ex \@minus .2ex}
{.1ex}
{\normalfont\small\sffamily\bfseries}}

%----------------------------------------------------------------------------------------
%	HYPERLINKS IN THE DOCUMENTS
%----------------------------------------------------------------------------------------

% For an unclear reason, the package should be loaded now and not later
\usepackage{hyperref}
\hypersetup{hidelinks,backref=true,pagebackref=true,hyperindex=true,colorlinks=false,breaklinks=true,urlcolor= mycolor,bookmarks=true,bookmarksopen=false,pdftitle={Title},pdfauthor={Author}}

%----------------------------------------------------------------------------------------
%	CHAPTER HEADINGS
%----------------------------------------------------------------------------------------

% The set-up below should be (sadly) manually adapted to the overall margin page septup controlled by the geometry package loaded in the main.tex document. It is possible to implement below the dimensions used in the goemetry package (top,bottom,left,right)... TO BE DONE

\newcommand{\thechapterimage}{}
\newcommand{\chapterimage}[1]{\renewcommand{\thechapterimage}{#1}}

% Numbered chapters with mini tableofcontents
\def\thechapter{\arabic{chapter}}
\def\@makechapterhead#1{
\thispagestyle{empty}
{\centering \normalfont\sffamily
\ifnum \c@secnumdepth >\m@ne
\if@mainmatter
\startcontents
\begin{tikzpicture}[remember picture,overlay]
\node at (current page.north west)
{\begin{tikzpicture}[remember picture,overlay]

% Comentar la linea para retirar la imagen
%\node[anchor=north west,inner sep=0pt] at (0,0) {\includegraphics[width=\paperwidth]{\thechapterimage}};

%%%%%%%%%%%%%%%%%%%%%%%%%%%%%%%%%%%%%%%%%%%%%%%%%%%%%%%%%%%%%%%%%%%%%%%%%%%%%%%%%%%%%
% Commenting the 3 lines below removes the small contents box in the chapter heading
%\fill[color=mycolor!10!white,opacity=.6] (10cm, -6.5cm) rectangle (22cm,-30cm);
%\node[anchor=north west] at (10cm,-6.5cm) {\parbox[t][8cm][t]{6.5cm}{\huge\bfseries\flushleft \printcontents{l}{1}{\setcounter{tocdepth}{2}}}};
\draw[anchor=west] (5cm,-5cm) node [rounded corners=20pt,fill=mycolor!10!white,text opacity=1,draw=mycolor,draw opacity=1,line width=1.5pt,fill opacity=.6,inner sep=12pt]{\huge\sffamily\bfseries\textcolor{black}{\thechapter. #1\strut\makebox[22cm]{}}}; %(5cm,-3cm) posicion del draw
%%%%%%%%%%%%%%%%%%%%%%%%%%%%%%%%%%%%%%%%%%%%%%%%%%%%%%%%%%%%%%%%%%%%%%%%%%%%%%%%%%%%%
\end{tikzpicture}};
\end{tikzpicture}}
\par\vspace*{150\p@} %Espacio para el empiezo de texto
\fi
\fi}

% Unnumbered chapters without mini tableofcontents (could be added though) %Index,Biblio, Contents, etc
\def\@makeschapterhead#1{
\thispagestyle{empty}
{\centering \normalfont\sffamily
\ifnum \c@secnumdepth >\m@ne
\if@mainmatter
\begin{tikzpicture}[remember picture,overlay]
\node at (current page.north west)
{\begin{tikzpicture}[remember picture,overlay]
%\node[anchor=north west,inner sep=0pt] at (0,0) {\includegraphics[width=\paperwidth]{\thechapterimage}}; %Remover la imagen 
\draw[anchor=west] (5cm,-5cm) node [rounded corners=20pt,fill=mycolor!10!white,fill opacity=.6,inner sep=12pt,text opacity=1,draw=mycolor,draw opacity=1,line width=1.5pt]{\huge\sffamily\bfseries\textcolor{black}{#1\strut\makebox[22cm]{}}};
\end{tikzpicture}};
\end{tikzpicture}}
\par\vspace*{180\p@} %Espacio para el empiezo de texto
\fi
\fi
}
\makeatother % Insert the commands.tex file which contains the majority of the structure behind the template


\begin{document}

%Cambiando al español :p o a lo que queres

%\renewcommand*\contentsname{Contenidos} %Para cambiar el nombre por defecto'Contents' por Contenidos aqui por alguna razon no cambia esta en content
\renewcommand*{\lstlistingname}{Codigo}
\renewcommand*{\lstlistlistingname}{Lista de \lstlistingname s} 
\renewcommand{\chaptername}{Capítulo}% Para cambiar el nombre por defecto'Chapter' por Capitulo
\renewcommand*\contentsname{Contenidos} %Para cambiar el nombre por defecto'Contents' por Contenidos
%----------------------------------------------------------------------------------------
%	TITLE PAGE
%----------------------------------------------------------------------------------------

\begingroup
\thispagestyle{empty}
\AddToShipoutPicture*{\put(6,5){\includegraphics[scale=1]{background}}} % Image background
\centering
\vspace*{9cm}
\par\normalfont\fontsize{35}{35}\sffamily\selectfont
Introducción a los algoritmos en c++\par % Book title
\vspace*{1cm}
{\Huge OscarGauss :P}\par % Author name
\endgroup

%----------------------------------------------------------------------------------------
%	COPYRIGHT PAGE
%----------------------------------------------------------------------------------------

\newpage
~\vfill
\thispagestyle{empty}

\noindent Copyright \copyright\ 2013 John Smith\\ % Copyright notice

\noindent \textsc{Published by Publisher}\\ % Publisher

\noindent \textsc{book-website.com}\\ % URL

\noindent Licensed under the Creative Commons Attribution-NonCommercial 3.0 Unported License (the ``License''). You may not use this file except in compliance with the License. You may obtain a copy of the License at \url{http://creativecommons.org/licenses/by-nc/3.0}. Unless required by applicable law or agreed to in writing, software distributed under the License is distributed on an \textsc{``as is'' basis, without warranties or conditions of any kind}, either express or implied. See the License for the specific language governing permissions and limitations under the License.\\ % License information

\noindent \textit{First printing, March 2013} % Printing/edition date

%----------------------------------------------------------------------------------------
%	TABLE OF CONTENTS
%----------------------------------------------------------------------------------------

%\chapterimage{chapter_head_000.pdf} % Table of contents heading image

\pagestyle{empty} % No headers

\tableofcontents % Print the table of contents itself

\cleardoublepage % Forces the first chapter to start on an odd page so it's on the right

\pagestyle{fancy} % Print headers again

%----------------------------------------------------------------------------------------
%	CHAPTERS
%----------------------------------------------------------------------------------------
%----------------------------------------------------------------------------------------
%	CHAPTER 1
%----------------------------------------------------------------------------------------

%\chapterimage{chapter_head_1.pdf} % Chapter heading image

\chapter{Introducción a la programación en c++}
\renewcommand{\chaptername}{Lecture}
\section{Introducción}\index{Introducción}

\subsection{¿Qué es un Lenguaje de Programación?}\index{Introducción!¿Qué es un Lenguaje de Programación?}

Antes de hablar de C++, es necesario explicar que un lenguaje de programación es una herramienta que nos permite comunicarnos e instruir a la computadora para que realice una tarea específica. Cada lenguaje de programación posee una sintaxis y un léxico particular, es decir, forma de escribirse que es diferente en cada uno por la forma que fue creado y por la forma que trabaja su compilador para revisar, acomodar y reservar el mismo programa en memoria.

Existen muchos lenguajes de programación de entre los que se destacan los siguientes:

\begin{itemize}
\item C
\item C++
\item Basic
\item Ada
\item Java
\item Pascal
\item Python
\item Fortran
\item Smalltalk
\end{itemize}

\subsection{Historia de C++}\index{Introducción!Historia de C++}

C++ es un lenguaje de programación creado por Bjarne Stroustrup en los laboratorios de At\&T en 1983. Stroustrup tomó como base el lenguaje de programación más popular en aquella época el cual era C.

El C++ es un derivado del mítico lenguaje C, el cual fue creado en la década de los 70 por la mano del finado Dennis Ritchie para la programación del sistema operativo (un sistema parecido a Unix es GNU/Linux), el cual surgió como un lenguaje orientado a la programación de sistemas (System Programming) y de herramientas (Utilities) recomendado sobre todo para programadores expertos, y que no llevaba implementadas muchas funciones que hacen a un lenguaje más comprensible.

Sin embargo, aunque esto en un inicio se puede convertir en un problema, en la práctica es su mayor virtud, ya que permite al programador un mayor control sobre lo que está haciendo. Años más tarde, un programador llamado Bjarne Stroustrup, creo lo que se conoce como C++.

Necesitaba ciertas facilidades de programación, incluidas en otros lenguajes pero que C no soportaba, al menos directamente, como son las llamadas clases y objetos, principios usados en la programación actual. Para ello rediseñó C, ampliando sus posibilidades pero manteniendo su mayor cualidad, la de permitir al programador en todo momento tener controlado lo que está haciendo, consiguiendo así una mayor rapidez que no se conseguiría en otros lenguajes.

C++ pretende llevar a C a un nuevo paradigma de clases y objetos con los que se realiza una comprensión más humana basándose en la construcción de objetos, con características propias solo de ellos, agrupados en clases. Es decir, si yo quisiera hacer un programa sobre animales, crearía una clase llamada animales, en la cual cada animal, por ejemplo un pato, sería un objeto, de tal manera que se ve el intento de esta forma de programar por ser un fiel reflejo de cómo los humanos (en teoría) manejamos la realidad.

Se dice que nuestro cerebro trabaja de forma relacional (relacionando hechos), es por ello que cada vez que recuerdas algo, (cuentas un hecho), termina siendo diferente (se agregan u omiten partes).

\subsection{¿Qué es C++?}\index{Introducción!¿Qué es C++?}

C++ es un lenguaje de programación orientado a objetos que toma la base del lenguaje C y le agrega la capacidad de abstraer tipos como en Smalltalk.

La intención de su creación fue el extender al exitoso lenguaje de programación C con mecanismos que permitieran la manipulación de objetos. En ese sentido, desde el punto de vista de los lenguajes orientados a objetos, el C++ es un lenguaje híbrido.

Posteriormente se añadieron facilidades de programación genérica, que se sumó a los otros dos paradigmas que ya estaban admitidos (programación estructurada y la programación orientada a objetos). Por esto se suele decir que el C++ es un lenguaje de programación multiparadigma.

\subsection{Herramientas Necesarias}\index{Introducción!Herramientas Necesarias}

Las principales herramientas necesarias para escribir un programa en C++ son las siguientes:

\begin{enumerate}
\item Un equipo ejecutando un sistema operativo.
\item Un compilador de C++
	\begin{itemize}
	\item Windows MingW (GCC para Windows) o MSVC (compilador de microsoft con versión gratuita)
	\item Linux (u otros UNIX): g++
	\item Mac (con el compilador Xcode)
	\end{itemize}
\item Un editor cualquiera de texto, o mejor un entorno de desarrollo (IDE)
	\begin{itemize}
	\item Windows:
		\begin{itemize}
		\item Microsoft Visual C++ (conocido por sus siglas MSVC). Incluye compilador y posee una versión gratuita (versión express)
		\item Bloc de notas (no recomendado)
		\item Editor Notepad++
		\item DevCpp (incluye MingW - en desuso, no recomendado, incluye también un compilador)
		\item Code::Blocks
		\end{itemize}
	\item Linux (o re-compilación en UNIX):
		\begin{itemize}
		\item Gedit
		\item Kate
		\item KDevelop
		\item Code::Blocks
		\item SciTE
		\item GVim
		\end{itemize}
	\item Mac:
		\begin{itemize}
		\item Xcode (con el compilador trae una IDE para poder programar)
		\end{itemize}
	\end{itemize}
\item Tiempo para practicar
\item Paciencia
\end{enumerate}

\textbf{Adicional}
\begin{itemize}
\item Inglés (Recomendado)
\item Estar familiarizado con C u otro lenguaje derivado (PHP, Python, etc).
\end{itemize}

Es recomendable tener conocimientos de C, debido a que C++ es una mejora de C, tener los conocimientos sobre este te permitira avanzar mas rapido y comprender aun mas. Tambien, hay que recordar que C++, admite C, por lo que se puede programar (reutilizar), funciones de C que se puedan usar en C++.

Aunque No es obligacion aprender C, es recomendable tener nociones sobre la programación orientada a objetos en el caso de no tener conocimientos previos de programación estructurada. Asimismo, muchos programadores recomiendan no saber C para saber C++, por ser el primero de ellos un lenguaje imperativo o procedimental y el segundo un lenguaje de programación orientado a objetos.

\subsection{Consejos iniciales antes de programar}\index{Introducción!Consejos iniciales antes de programar}

Con la práctica, se puede observar que se puede confundir a otros programadores con el código que se haga. Antes de siquiera hacer una línea de código, si se trabaja con otros programadores, ha de tenerse en cuenta que todos deben escribir de una forma similar el código, para que de forma global puedan corregir el código en el caso de que hubieran errores o rastrearlos en el caso de haberlos.

También es muy recomendable hacer uso de comentarios (comenta todo lo que puedas, hay veces que lo que parece obvio para ti, no lo es para los demás) y tratar de hacer un código limpio y comprensible, especificando detalles y haciendo tabulaciones, aunque te tome un poco mas de tiempo, es posible que mas adelante lo agradezcas tu mismo.

\subsection{Ejemplos}\index{Introducción!Ejemplos}

\begin{lstlisting}[style=Cpp,label=ejemplo-c++,caption=Ejemplo de C++]
#include <iostream>

using namespace std;

int main()
{
    int numero;
    cin>>numero;
    if(numero%2==0){
        cout<<"El numero es par\n";
    }else{
        cout<<"EL numero es impar\n";
    }
    return 0;
}
\end{lstlisting}


\section{Lo mas basico}\index{Lo mas basico}

\subsection{Proceso de desarrollo de un programa}\index{Lo mas basico!Proceso de desarrollo de un programa}

Si se desea escribir un programa en C++ se debe ejecutar como mínimo los siguientes pasos:
\begin{enumerate}
\item Escribir con un editor de texto plano un programa sintácticamente válido o usar un entorno de desarrollo (IDE) apropiado para tal fin
\item Compilar el programa y asegurarse de que no han habido errores de compilación
\item Ejecutar el programa y comprobar que no hay errores de ejecución
\end{enumerate}
Este último paso es el más costoso, por que en programas grandes, averiguar si hay o no un fallo prácticamente puede ser una tarea totémica.

Un archivo de C++ tiene la extención \textit{cpp} a continuación se escribe el siguiente codigo en c++ del archivo 'hola.cpp'

\begin{lstlisting}[style=Cpp, label=hola-mundo, caption=Hola Mundo]
// Aquí generalmente se suele indicar qué se quiere con el programa a hacer
// Programa que muestra 'Hola mundo' por pantalla y finaliza

// Aquí se sitúan todas las librerias que se vayan a usar con include,
// que se verá posteriormente
#include <iostream> // Esta libreria permite mostrar y leer datos por consola

int main()
{
    // Este tipo de líneas de código que comienzan por '//' son comentarios
    // El compilador los omite, y sirven para ayudar a otros programadores o
    // a uno mismo en caso de volver a revisar el código
    // Es una práctica sana poner comentarios donde se necesiten,

    std::cout << "Hola Mundo" << std::endl;

    // Mostrar por std::cout el mensaje Hola Mundo y comienza una nueva línea
    
    return 0;

    // se devuelve un 0.
    //que en este caso quiere decir que la salida se ha efectuado con éxito.
}
\end{lstlisting}

Mediante simple inspección, el código parece enorme, pero el compilador lo único que leerá para la creación del programa es lo siguiente:
\begin{lstlisting}[style=Cpp, label=hola-mundo-compilador, caption=Hola Mundo Comp]
#include <iostream>
int main(void){ std::cout << "Hola Mundo" << std::endl; return 0; }
\end{lstlisting}

Como se puede observar, este código y el original no difieren en mucho salvo en los saltos de línea y que los comentarios, de los que se detallan posteriormente, están omitidos y tan sólo ha quedado "el esqueleto" del código legible para el compilador. Para el compilador, todo lo demás, sobra.

Aquí otro ejemplo

\begin{lstlisting}[style=Cpp, label=hola-mundo-compilador, caption=Hola Mundo Comp]
#include <iostream>

int main()
{
    std::cout << "Hola Mundo" << std::endl;
    std::cout << "Hello World" << std::endl;
    std::cout << "Hallo Welt" << std::endl;
    return 0;
}
\end{lstlisting}

Para hacer el código mas corto debemos incluir \textit{using namespace std} con lo que le estamos diciendo al compilador que usaremos el espacio de nombres std por lo que no tendremos que incluirlo cuando usemos elementos de este espacio de nombres, como pueden ser los objetos cout y cin, que representan el flujo de salida estándar (típicamente la pantalla o una ventana de texto) y el flujo de entrada estándar (típicamente el teclado).

Se veria de la siguiente manera:

\begin{lstlisting}[style=Cpp, label=hola-mundo-compilador, caption=Hola Mundo Comp]
#include <iostream>

using namespace std;

int main()
{
    cout << "Hola Mundo" << endl;
    cout << "Hello World" << endl;
    cout << "Hallo Welt" << endl;
    return 0;
}
\end{lstlisting}

Los pasos siguientes son para una compilación en GNU o sistema operativo Unix, para generar el ejecutable del programa se compila con g++ de la siguiente forma:

\begin{listing}[style=consola, numbers=none,label=some-code,caption=Some Code]
| g++ hola.cpp -o hola
\end{listing}

Para poder ver los resultados del programa en acción, se ejecuta el programa de la siguiente forma:

\begin{listing}[style=consola, numbers=none,label=some-code,caption=Some Code]
|./hola
\end{listing}

Y a continuación se debe mostrar algo como lo siguiente:

\begin{listing}[style=consola, numbers=none,label=some-code,caption=Some Code]
Hola Mundo
\end{listing}

\subsection{Comentarios}\index{Lo mas basico!Comentarios}
\subsection{Impresión}\index{Lo mas basico!Impresión}
\subsection{Lectura}\index{Lo mas basico!Lectura}
\subsection{Lectura}\index{Lo mas basico!Lectura}

\begin{listing}[style=consola, numbers=none,label=some-code,caption=Some Code]
$ gcc  -o hello hello.c
\end{listing}

%----------------------------------------------------------------------------------------
%	CHAPTER 1
%----------------------------------------------------------------------------------------

%\chapterimage{chapter_head_1.pdf} % Chapter heading image

\chapter{Introducción a la programación en c++}
\renewcommand{\chaptername}{Lecture}
\section{Introducción}\index{Introducción}

\subsection{¿Qué es un Lenguaje de Programación?}\index{Introducción!¿Qué es un Lenguaje de Programación?}

Antes de hablar de C++, es necesario explicar que un lenguaje de programación es una herramienta que nos permite comunicarnos e instruir a la computadora para que realice una tarea específica. Cada lenguaje de programación posee una sintaxis y un léxico particular, es decir, forma de escribirse que es diferente en cada uno por la forma que fue creado y por la forma que trabaja su compilador para revisar, acomodar y reservar el mismo programa en memoria.

Existen muchos lenguajes de programación de entre los que se destacan los siguientes:

\begin{itemize}
\item C
\item C++
\item Basic
\item Ada
\item Java
\item Pascal
\item Python
\item Fortran
\item Smalltalk
\end{itemize}

\subsection{Historia de C++}\index{Introducción!Historia de C++}

C++ es un lenguaje de programación creado por Bjarne Stroustrup en los laboratorios de At\&T en 1983. Stroustrup tomó como base el lenguaje de programación más popular en aquella época el cual era C.

El C++ es un derivado del mítico lenguaje C, el cual fue creado en la década de los 70 por la mano del finado Dennis Ritchie para la programación del sistema operativo (un sistema parecido a Unix es GNU/Linux), el cual surgió como un lenguaje orientado a la programación de sistemas (System Programming) y de herramientas (Utilities) recomendado sobre todo para programadores expertos, y que no llevaba implementadas muchas funciones que hacen a un lenguaje más comprensible.

Sin embargo, aunque esto en un inicio se puede convertir en un problema, en la práctica es su mayor virtud, ya que permite al programador un mayor control sobre lo que está haciendo. Años más tarde, un programador llamado Bjarne Stroustrup, creo lo que se conoce como C++.

Necesitaba ciertas facilidades de programación, incluidas en otros lenguajes pero que C no soportaba, al menos directamente, como son las llamadas clases y objetos, principios usados en la programación actual. Para ello rediseñó C, ampliando sus posibilidades pero manteniendo su mayor cualidad, la de permitir al programador en todo momento tener controlado lo que está haciendo, consiguiendo así una mayor rapidez que no se conseguiría en otros lenguajes.

C++ pretende llevar a C a un nuevo paradigma de clases y objetos con los que se realiza una comprensión más humana basándose en la construcción de objetos, con características propias solo de ellos, agrupados en clases. Es decir, si yo quisiera hacer un programa sobre animales, crearía una clase llamada animales, en la cual cada animal, por ejemplo un pato, sería un objeto, de tal manera que se ve el intento de esta forma de programar por ser un fiel reflejo de cómo los humanos (en teoría) manejamos la realidad.

Se dice que nuestro cerebro trabaja de forma relacional (relacionando hechos), es por ello que cada vez que recuerdas algo, (cuentas un hecho), termina siendo diferente (se agregan u omiten partes).

\subsection{¿Qué es C++?}\index{Introducción!¿Qué es C++?}

C++ es un lenguaje de programación orientado a objetos que toma la base del lenguaje C y le agrega la capacidad de abstraer tipos como en Smalltalk.

La intención de su creación fue el extender al exitoso lenguaje de programación C con mecanismos que permitieran la manipulación de objetos. En ese sentido, desde el punto de vista de los lenguajes orientados a objetos, el C++ es un lenguaje híbrido.

Posteriormente se añadieron facilidades de programación genérica, que se sumó a los otros dos paradigmas que ya estaban admitidos (programación estructurada y la programación orientada a objetos). Por esto se suele decir que el C++ es un lenguaje de programación multiparadigma.

\subsection{Herramientas Necesarias}\index{Introducción!Herramientas Necesarias}

Las principales herramientas necesarias para escribir un programa en C++ son las siguientes:

\begin{enumerate}
\item Un equipo ejecutando un sistema operativo.
\item Un compilador de C++
	\begin{itemize}
	\item Windows MingW (GCC para Windows) o MSVC (compilador de microsoft con versión gratuita)
	\item Linux (u otros UNIX): g++
	\item Mac (con el compilador Xcode)
	\end{itemize}
\item Un editor cualquiera de texto, o mejor un entorno de desarrollo (IDE)
	\begin{itemize}
	\item Windows:
		\begin{itemize}
		\item Microsoft Visual C++ (conocido por sus siglas MSVC). Incluye compilador y posee una versión gratuita (versión express)
		\item Bloc de notas (no recomendado)
		\item Editor Notepad++
		\item DevCpp (incluye MingW - en desuso, no recomendado, incluye también un compilador)
		\item Code::Blocks
		\end{itemize}
	\item Linux (o re-compilación en UNIX):
		\begin{itemize}
		\item Gedit
		\item Kate
		\item KDevelop
		\item Code::Blocks
		\item SciTE
		\item GVim
		\end{itemize}
	\item Mac:
		\begin{itemize}
		\item Xcode (con el compilador trae una IDE para poder programar)
		\end{itemize}
	\end{itemize}
\item Tiempo para practicar
\item Paciencia
\end{enumerate}

\textbf{Adicional}
\begin{itemize}
\item Inglés (Recomendado)
\item Estar familiarizado con C u otro lenguaje derivado (PHP, Python, etc).
\end{itemize}

Es recomendable tener conocimientos de C, debido a que C++ es una mejora de C, tener los conocimientos sobre este te permitira avanzar mas rapido y comprender aun mas. Tambien, hay que recordar que C++, admite C, por lo que se puede programar (reutilizar), funciones de C que se puedan usar en C++.

Aunque No es obligacion aprender C, es recomendable tener nociones sobre la programación orientada a objetos en el caso de no tener conocimientos previos de programación estructurada. Asimismo, muchos programadores recomiendan no saber C para saber C++, por ser el primero de ellos un lenguaje imperativo o procedimental y el segundo un lenguaje de programación orientado a objetos.

\subsection{Consejos iniciales antes de programar}\index{Introducción!Consejos iniciales antes de programar}

Con la práctica, se puede observar que se puede confundir a otros programadores con el código que se haga. Antes de siquiera hacer una línea de código, si se trabaja con otros programadores, ha de tenerse en cuenta que todos deben escribir de una forma similar el código, para que de forma global puedan corregir el código en el caso de que hubieran errores o rastrearlos en el caso de haberlos.

También es muy recomendable hacer uso de comentarios (comenta todo lo que puedas, hay veces que lo que parece obvio para ti, no lo es para los demás) y tratar de hacer un código limpio y comprensible, especificando detalles y haciendo tabulaciones, aunque te tome un poco mas de tiempo, es posible que mas adelante lo agradezcas tu mismo.

\subsection{Ejemplos}\index{Introducción!Ejemplos}

\begin{lstlisting}[style=Cpp,label=ejemplo-c++,caption=Ejemplo de C++]
#include <iostream>

using namespace std;

int main()
{
    int numero;
    cin>>numero;
    if(numero%2==0){
        cout<<"El numero es par\n";
    }else{
        cout<<"EL numero es impar\n";
    }
    return 0;
}
\end{lstlisting}


\section{Lo mas basico}\index{Lo mas basico}

\subsection{Proceso de desarrollo de un programa}\index{Lo mas basico!Proceso de desarrollo de un programa}

Si se desea escribir un programa en C++ se debe ejecutar como mínimo los siguientes pasos:
\begin{enumerate}
\item Escribir con un editor de texto plano un programa sintácticamente válido o usar un entorno de desarrollo (IDE) apropiado para tal fin
\item Compilar el programa y asegurarse de que no han habido errores de compilación
\item Ejecutar el programa y comprobar que no hay errores de ejecución
\end{enumerate}
Este último paso es el más costoso, por que en programas grandes, averiguar si hay o no un fallo prácticamente puede ser una tarea totémica.

Un archivo de C++ tiene la extención \textit{cpp} a continuación se escribe el siguiente codigo en c++ del archivo 'hola.cpp'

\begin{lstlisting}[style=Cpp, label=hola-mundo, caption=Hola Mundo]
// Aquí generalmente se suele indicar qué se quiere con el programa a hacer
// Programa que muestra 'Hola mundo' por pantalla y finaliza

// Aquí se sitúan todas las librerias que se vayan a usar con include,
// que se verá posteriormente
#include <iostream> // Esta libreria permite mostrar y leer datos por consola

int main()
{
    // Este tipo de líneas de código que comienzan por '//' son comentarios
    // El compilador los omite, y sirven para ayudar a otros programadores o
    // a uno mismo en caso de volver a revisar el código
    // Es una práctica sana poner comentarios donde se necesiten,

    std::cout << "Hola Mundo" << std::endl;

    // Mostrar por std::cout el mensaje Hola Mundo y comienza una nueva línea
    
    return 0;

    // se devuelve un 0.
    //que en este caso quiere decir que la salida se ha efectuado con éxito.
}
\end{lstlisting}

Mediante simple inspección, el código parece enorme, pero el compilador lo único que leerá para la creación del programa es lo siguiente:
\begin{lstlisting}[style=Cpp, label=hola-mundo-compilador, caption=Hola Mundo Comp]
#include <iostream>
int main(void){ std::cout << "Hola Mundo" << std::endl; return 0; }
\end{lstlisting}

Como se puede observar, este código y el original no difieren en mucho salvo en los saltos de línea y que los comentarios, de los que se detallan posteriormente, están omitidos y tan sólo ha quedado "el esqueleto" del código legible para el compilador. Para el compilador, todo lo demás, sobra.

Aquí otro ejemplo

\begin{lstlisting}[style=Cpp, label=hola-mundo-compilador, caption=Hola Mundo Comp]
#include <iostream>

int main()
{
    std::cout << "Hola Mundo" << std::endl;
    std::cout << "Hello World" << std::endl;
    std::cout << "Hallo Welt" << std::endl;
    return 0;
}
\end{lstlisting}

Para hacer el código mas corto debemos incluir \textit{using namespace std} con lo que le estamos diciendo al compilador que usaremos el espacio de nombres std por lo que no tendremos que incluirlo cuando usemos elementos de este espacio de nombres, como pueden ser los objetos cout y cin, que representan el flujo de salida estándar (típicamente la pantalla o una ventana de texto) y el flujo de entrada estándar (típicamente el teclado).

Se veria de la siguiente manera:

\begin{lstlisting}[style=Cpp, label=hola-mundo-compilador, caption=Hola Mundo Comp]
#include <iostream>

using namespace std;

int main()
{
    cout << "Hola Mundo" << endl;
    cout << "Hello World" << endl;
    cout << "Hallo Welt" << endl;
    return 0;
}
\end{lstlisting}

Los pasos siguientes son para una compilación en GNU o sistema operativo Unix, para generar el ejecutable del programa se compila con g++ de la siguiente forma:

\begin{listing}[style=consola, numbers=none,label=some-code,caption=Some Code]
| g++ hola.cpp -o hola
\end{listing}

Para poder ver los resultados del programa en acción, se ejecuta el programa de la siguiente forma:

\begin{listing}[style=consola, numbers=none,label=some-code,caption=Some Code]
|./hola
\end{listing}

Y a continuación se debe mostrar algo como lo siguiente:

\begin{listing}[style=consola, numbers=none,label=some-code,caption=Some Code]
Hola Mundo
\end{listing}

\subsection{Comentarios}\index{Lo mas basico!Comentarios}
\subsection{Impresión}\index{Lo mas basico!Impresión}
\subsection{Lectura}\index{Lo mas basico!Lectura}
\subsection{Lectura}\index{Lo mas basico!Lectura}

\begin{listing}[style=consola, numbers=none,label=some-code,caption=Some Code]
$ gcc  -o hello hello.c
\end{listing}

%----------------------------------------------------------------------------------------
%	BIBLIOGRAPHY
%----------------------------------------------------------------------------------------

\chapter*{Bibliography}
\addcontentsline{toc}{chapter}{\textcolor{mycolor}{Bibliography}}
\section*{Books}
\addcontentsline{toc}{section}{Books}
\printbibliography[heading=bibempty,type=book]
\section*{Articles}
\addcontentsline{toc}{section}{Articles}
\printbibliography[heading=bibempty,type=article]

%----------------------------------------------------------------------------------------
%	INDEX
%----------------------------------------------------------------------------------------

\cleardoublepage
\phantomsection
\setlength{\columnsep}{0.75cm}
\addcontentsline{toc}{chapter}{\textcolor{mycolor}{Index}}
\printindex

%----------------------------------------------------------------------------------------

%Listar las estrofas de códigos
\lstlistoflistings

\end{document}